\chapter{Fundamentos}
\label{cha:fundamentos}
	
	\lipsum

\section{Section 1}
\label{sec:industria4}

	\lipsum

	\begin{figure}[htb]
		\centering
		\includegraphics[width=1\textwidth]{sample.png}
		\caption{Sample image.}
		\label{fig:sample2}
		\fonte{\citeonline{wahlster2013industrie} (adaptado). }
	\end{figure}

	Conforme listados na \autoref{tab:tabela1}.
	
	\begin{table}[htb]
		\centering
		\footnotesize
		\begin{tabular}{p{3cm}p{12cm}}
			\hline
			\textbf{Princípio} & \textbf{Descrição} \\
			
			\hline
			Interoperabilidade &
			Capacidade das coisas (máquinas, dispositivos, sensores, pessoas, etc) de comunicarem entre si dentro de um sistema por meio de padrões definidos. \\
			
			\hline
			Transparência de informação &
			Tornar acessíveis informações úteis para os demais dispositivos conectados à rede. Informações do mundo virtual como documentos eletrônicos, desenhos, modelos de simulação; e informações sobre o mundo real, como posição, dados de sensores de temperatura, vibração, etc. \\
			
			\hline
			Descentralização de decisões &
			Permitir a tomada de decisões baseada nas informações coletados pelo próprio dispositivo e dar ao dispositivo autonomia para decidir qual será sua próxima função/operação. Desta forma, um planejamento ou controle central de processos produtivos não se faz essencial e o sistema de produção se torna menos hierarquizado. \\
			
			\hline
			Assistência técnica &
			Devido à complexidade da produção, com redes complexas e tomada decisões descentralizadas, os seres humanos precisam ser auxiliados por sistemas de assistência de forma a dar compreensibilidade ao processo e às tomadas de decisão necessárias. Os sistemas de assistência devem agregar e tornar visualizáveis as informações de maneira compreensível. \\
			
			\hline
		\end{tabular}
		\caption{Tabela 1.}
		\label{tab:tabela1}
	\end{table}

	\subsection{Subsection }
	\label{sub:rami4}
	
	\lipsum